\clearpage\fancyhead[L]{\small JUNO.}\SourceLeaf{397}
\begin{center}\textbf{GAUSS TO BODE.}\par\vspace{4pt}\hrule\vspace{4pt}\fontsize{10}{12}\selectfont Astronomisches Jahrbuch für 1819. pp. 219..221. Berlin 1816.\par\vspace{4pt}\hrule\vspace{4pt}\end{center}
\begin{flushright}Göttingen, 10 August 1816.\end{flushright}

\par\indent At your request, my esteemed friend, I am pleased to send you the enclosed ephemerides for the next apparitions of Pallas and Juno, with a few remarks.

\par\indent The ephemeris for Pallas was calculated by Mr. \textsc{Westphal} of Schwerin, who is now devoting himself here to astronomy with great zeal. As has always been done for the past 5 years, the perturbations have been taken into account according to my theory, so the closest agreement may be expected. Mr. \textsc{Westphal} calculated the next opposition in advance as follows:
\begin{center}\fontsize{10}{12}\selectfont
11 July 1817. $8^{\mathrm u}\;45^{\mathrm m}\;5^{\mathrm s}$ Mean time in Göttingen\par
\setlength{\tabcolsep}{3pt}\begin{tabular}{lcr}
True longitude & $.\;.\;.\;.\;.\;.\;.\;.\;.\;.\;.\;.\;.\;.$ & $289^\circ\;3'\;22''$\\
Geocentric latitude & $.\;.\;.\;.\;.\;.\;.\;.\;.\;.\;.$ & $43\;33\;1$\\
\end{tabular}\end{center}

\par\indent Pallas will be very faint next year, at only magnitude 10–11. Here is a survey of its brightness at the last oppositions and the coming one:
\begin{center}\fontsize{10}{12}\selectfont\setlength{\tabcolsep}{11pt}\begin{tabular}{lr|lr}
1812 & $0.01666$ & 1816 & $0.05997$\\
1813 & $0.01475$ & 1817 & $0.01289$\\
1814 & $0.05476$ & 1818 & $0.02122$\\
\end{tabular}\end{center}

\Needspace{9\baselineskip}\par\indent I myself observed Juno twice with the ring micrometer:
\begin{center}\fontsize{10}{12}\selectfont\setlength{\tabcolsep}{5pt}\begin{tabular}{l|c|c|c}
1816 & Mean time & Right ascension & \shortstack{Southern\\declination}\\\hline
June 12 & $10^{\mathrm u}\;47^{\mathrm m}\;27^{\mathrm s}5$ & $251^\circ\;46'\;8''0$ & $3^\circ\;42'\;4''7$\\
13 & $10\;27\;34.7$ & $251\;33\;47.6$ & $3\;41\;11.7$\\
\end{tabular}\end{center}

\Needspace{11\baselineskip}\par\indent With the unusually bad weather last spring, no further observations could be made. I derived the following results for the opposition from these observations:
\begin{center}\fontsize{10}{12}\selectfont
3 June 1816. $4^{\mathrm u}\;20^{\mathrm m}\;20^{\mathrm s}$ Mean time in Göttingen\par
\setlength{\tabcolsep}{3pt}\begin{tabular}{lcr}
True longitude & $.\;.\;.\;.\;.\;.\;.\;.\;.\;.\;.\;.$ & $252^\circ\;51'\;39''2$\\
Geocentric latitude & $.\;.\;.\;.\;.\;.\;.\;.\;.$ & $18\;34\;3.1$ N.\\
\end{tabular}\end{center}

\Needspace{9\baselineskip}\par\indent Mr. \textsc{Enke}, who later repeated this calculation and also used the Vilna observations, found
\begin{center}\fontsize{10}{12}\selectfont
3 June \quad $4^{\mathrm u}\;17^{\mathrm m}\;31^{\mathrm s}$ \quad Mean time in Göttingen\par
$252^\circ\;51'\;32''6$\par
$18\;34\;14.6$
\end{center}

\Needspace{15\baselineskip}\par\indent By combining the last four oppositions, I derived the following elements:
\SourceLeaf{398}\begin{center}\fontsize{9}{11}\selectfont\setlength{\tabcolsep}{2pt}\renewcommand{\arraystretch}{1.15}\begin{tabular}{lcr}
Mean longitude 6 Sept. 1817. $0^{\mathrm u}$ in Göttingen & $.\;.$ & $8^\circ\;51'\;41''$\\
Perihelion & $.\;.\;.\;.\;.\;.\;.\;.\;.\;.\;.\;.\;.\;.$ & $53\;16\;46$\\
Node & $.\;.\;.\;.\;.\;.\;.\;.\;.\;.\;.\;.\;.\;.\;.\;.$ & $171\;9\;41$\\
Inclination & $.\;.\;.\;.\;.\;.\;.\;.\;.\;.\;.\;.\;.\;.\;.$ & $13\;3\;49$\\
Eccentricity angle & $.\;.\;.\;.\;.\;.\;.\;.\;.\;.\;.$ & $14\;46\;0$\\
Mean daily tropical motion & $.\;.\;.\;.$ & $812''40187$\\
Logarithm of the semimajor axis & $.\;.\;.\;.\;.\;.$ & $0.4269223$\\
\end{tabular}\end{center}

\par\indent (It is probably superfluous to point out that these elements were calculated solely for use in computing the next ephemeris. As long as the perturbations have yet to be worked out, this method is the most convenient and the most accurate.)
\par\indent From these elements Prof. \textsc{Harding} calculated an ephemeris for the next apparition, and I enclose a copy for you. Mr. \textsc{Posselt} also carried out the same calculation. His ephemeris differs from this one throughout only in the seconds and will be printed in Mr. \textsc{von Lindenau}’s journal. The next opposition occurs as follows according to Prof. \textsc{Harding}’s calculation from the above elements.
\begin{center}\fontsize{10}{12}\selectfont
6 Sept. 1817. $4^{\mathrm u}\;10^{\mathrm m}\;30^{\mathrm s}$ mean time in Göttingen.\par
\setlength{\tabcolsep}{3pt}\begin{tabular}{lcr}
True longitude & $.\;.\;.\;.\;.\;.\;.\;.\;.\;.\;.\;.\;.$ & $343^\circ\;37'\;36''$\\
Geocentric latitude & $.\;.\;.\;.\;.\;.\;.\;.\;.\;.\;.$ & $3\;6\;9$ N.\\
Brightness in 1816 & $.\;.\;.\;.\;.\;.\;.\;.\;.\;.\;.$ & $0.01747$\\
\hfill1817 & $.\;.\;.\;.\;.\;.\;.\;.\;.\;.\;.$ & $0.11398$\\
\end{tabular}\end{center}

\par\indent Allow me to add a few remarks on variable stars. In the yearbook for 1818, Mr. \textsc{Koch} assigns a period of 295 days to the star in Virgo discovered by Prof. \textsc{Harding}, without saying what he bases it on. But this period does not agree with Prof. \textsc{Harding}’s observations. — In 1811 Prof. \textsc{Harding} noticed a second variable star in Aquarius, south of $\omega^2$. It fades gradually from the sixth magnitude until it becomes entirely invisible, in about $2\frac12$ months. Its brightness changes do not seem to follow any regular period at all. Its mean position in 1800 was
\begin{center}\fontsize{10}{12}\selectfont\setlength{\tabcolsep}{3pt}\begin{tabular}{lcr}
Right ascension & $.\;.\;.\;.\;.\;.\;.\;.\;.\;.\;.$ & $353^\circ\;21'\;48''5$\\
Annual. var. & $.\;.\;.\;.\;.\;.\;.\;.\;.\;.\;.\;.\;.\;.\;.$ & $+46.56$\\
Southern declination & $.\;.\;.\;.\;.\;.\;.\;.\;.\;.$ & $16^\circ\;23'\;13''4$\\
Annual decrease & $.\;.\;.\;.\;.\;.\;.\;.\;.\;.\;.\;.\;.\;.$ & $19.88$.\\
\end{tabular}\end{center}
